% =====================================================================
%  CARNET D'ENTRAINEMENT — FICHE 6 : Liaisons et structures de Lewis
%  THÈME 3 — Charge formelle et choix de structure
%  Fiche TUTO (à lire avant les exercices)
%  Compilation : pdflatex
% =====================================================================
\documentclass[11pt,a4paper]{article}
\usepackage[utf8]{inputenc}
\usepackage[T1]{fontenc}
\usepackage{lmodern}
\usepackage[french]{babel}
\usepackage{amsmath,amssymb,mathtools}
\usepackage{icomma}
\usepackage[version=4]{mhchem}
\usepackage{siunitx}
\sisetup{output-decimal-marker={,}}
\usepackage{xcolor}
\usepackage{colortbl}
\usepackage{tikz}
\usepackage[most]{tcolorbox}
\usepackage{enumitem}
\usepackage{array}
\usepackage{booktabs}
\usepackage{multicol}
\usepackage{fancyhdr}
\usepackage[hidelinks]{hyperref}
\usetikzlibrary{arrows.meta,positioning,calc,shapes.geometric,shapes.misc,fit,backgrounds}
\usepackage[a4paper,margin=2cm,headheight=26pt,headsep=8pt]{geometry}

% -------- palette --------
\definecolor{primary}{RGB}{146,95,160}
\definecolor{primarydark}{RGB}{92,52,104}
\definecolor{defcol}{RGB}{0,114,178}
\definecolor{thmcol}{RGB}{0,140,120}
\definecolor{excol}{RGB}{0,135,90}
\definecolor{methcol}{RGB}{90,90,100}
\definecolor{attncol}{RGB}{200,40,55}
\definecolor{astucecol}{RGB}{198,93,9}
\definecolor{atomCl}{RGB}{70,170,80}
\definecolor{atomO}{RGB}{220,55,45}
\definecolor{atomN}{RGB}{48,80,200}
\definecolor{atomH}{RGB}{235,235,235}
\definecolor{lonepair}{RGB}{120,94,200}
\definecolor{bondcol}{RGB}{60,60,70}
\definecolor{piecol}{RGB}{198,93,9}

% -------- boîtes --------
\tcbset{boxbase/.style={enhanced,breakable,boxrule=0.9pt,arc=2.5pt,
   left=3mm,right=3mm,top=2.2mm,bottom=2.2mm,
   fonttitle=\bfseries,coltitle=white,toptitle=0.6mm,bottomtitle=0.6mm}}
\newtcolorbox{defbox}[1][]{boxbase,colback=defcol!6,colframe=defcol,
   title={\textsc{À retenir}\ifx\relax#1\relax\relax\else\textmd{ : \itshape #1}\fi}}
\newtcolorbox{methbox}[1][]{boxbase,colback=methcol!8,colframe=methcol,sharp corners,
   title={\textsc{Méthode}\ifx\relax#1\relax\relax\else\textmd{ : \itshape #1}\fi}}
\newtcolorbox{exbox}[1][]{boxbase,colback=excol!6,colframe=excol,
   title={\textsc{Exemple}\ifx\relax#1\relax\relax\else\textmd{ : \itshape #1}\fi}}
\newtcolorbox{attnbox}[1][]{boxbase,colback=attncol!6,colframe=attncol,
   title={\textsc{Piège à éviter}\ifx\relax#1\relax\relax\else\textmd{ : \itshape #1}\fi}}
\newtcolorbox{astucebox}[1][]{boxbase,colback=astucecol!8,colframe=astucecol,
   title={\textsc{Astuce}\ifx\relax#1\relax\relax\else\textmd{ : \itshape #1}\fi}}

\newcommand{\CF}{\ensuremath{q_{\mathrm{f}}}}

% -------- en-tête --------
\pagestyle{fancy}\fancyhf{}
\fancyhead[L]{\small\textcolor{primarydark}{\textbf{Fiche 6 · Thème 3} — Charge formelle et choix de structure}}
\fancyhead[R]{\small\textcolor{primarydark}{\textsc{Tuto}}}
\fancyfoot[C]{\small\thepage}
\renewcommand{\headrulewidth}{0.8pt}
\renewcommand{\headrule}{\hbox to\headwidth{\color{primary}\leaders\hrule height \headrulewidth\hfill}}

\setlength{\parindent}{0pt}
\setlength{\parskip}{3pt}

\begin{document}

\begin{tcolorbox}[boxbase,colback=primary!12,colframe=primary,
   title={\Large Thème 3 — Charge formelle et choix de structure}]
\textbf{Objectif du thème.} Savoir \emph{calculer} la charge formelle (\CF) d'un atome dans une
structure de Lewis, \emph{vérifier} qu'elle est cohérente (la somme redonne la charge de
l'édifice), et surtout \emph{s'en servir} pour \textbf{départager} plusieurs structures possibles
et choisir la plus représentative.
\end{tcolorbox}

\section*{1. La charge formelle : une comptabilité d'électrons}

Une fois une structure de Lewis dessinée, on attribue à chaque atome une \textbf{charge formelle}.
Ce n'est pas une charge réelle : c'est un \emph{outil de comptabilité} qui dit si l'atome a, dans
cette structure, \emph{plus} ou \emph{moins} d'électrons que lorsqu'il est neutre et isolé.

\begin{defbox}[la formule à connaître par c\oe ur]
\[ \boxed{\;\CF = N_v - 2\times(\text{nb de doublets non liants}) - (\text{nb de liaisons})\;} \]
où $N_v$ est le nombre d'électrons de valence de l'atome neutre isolé (son numéro de colonne).
On compte les \textbf{doublets libres en entier} ($\times 2$ électrons), et chaque \textbf{liaison
apporte un seul électron propre} à l'atome (l'autre est attribué au partenaire).\\[2pt]
\footnotesize\emph{Valences utiles :} \ce{H}=1, \ce{C}=4, \ce{N}=5, \ce{O}=6, \ce{S}=6, \ce{Cl}=7.
\end{defbox}

\begin{defbox}[trois propriétés indispensables]
\begin{itemize}[leftmargin=1.5em,topsep=1pt,itemsep=2pt]
  \item \textbf{Somme = charge globale.} La somme des \CF{} de tous les atomes est égale à la charge
        de l'édifice ($0$ pour une molécule neutre, $q$ pour un ion). C'est le test de cohérence.
  \item \textbf{Meilleure structure = \CF{} minimales.} La structure la plus représentative est
        celle dont les charges formelles sont les plus \textbf{proches de zéro}.
  \item \textbf{Bon placement.} Les charges \emph{négatives} doivent se trouver sur les atomes les
        plus \textbf{électronégatifs} (souvent \ce{O}, \ce{F}\ldots).
\end{itemize}
\end{defbox}

\section*{2. Calculer une charge formelle : exemple guidé}

\begin{methbox}[la marche à suivre pour un atome]
\begin{enumerate}[leftmargin=1.7em,topsep=1pt,itemsep=1pt]
  \item Lis sa valence $N_v$ (numéro de colonne).
  \item Compte ses \textbf{doublets non liants} (paires libres) et multiplie par $2$.
  \item Compte ses \textbf{liaisons} — \emph{attention : une double liaison compte pour 2 liaisons}.
  \item Applique $\CF = N_v - 2\times(\text{doublets libres}) - (\text{liaisons})$.
\end{enumerate}
\end{methbox}

\begin{exbox}[les charges formelles de l'ion carbonate \ce{CO3^2-}]
Structure : un carbone central lié à trois oxygènes — \emph{un} \ce{O} par une double liaison,
\emph{deux} \ce{O} par une simple liaison.

\smallskip
\textbf{Carbone central :} $0$ doublet libre, $4$ liaisons (une double $+$ deux simples).\;
$\CF = 4 - 2\times 0 - 4 = \mathbf{0}.$

\smallskip
\textbf{Oxygène doublement lié :} $2$ doublets libres, $2$ liaisons.\;
$\CF = 6 - 2\times 2 - 2 = 6 - 4 - 2 = \mathbf{0}.$

\smallskip
\textbf{Chaque oxygène simplement lié :} $3$ doublets libres, $1$ liaison.\;
$\CF = 6 - 2\times 3 - 1 = 6 - 6 - 1 = \mathbf{-1}.$

\smallskip
\textbf{Bilan.}\quad $\CF$ totale $= 0 + 0 + (-1) + (-1) = -2$ : on retrouve la charge de l'ion, et
les charges $-$ sont portées par les oxygènes (les plus électronégatifs). Structure cohérente.
\end{exbox}

\section*{3. Choisir la meilleure structure}

Quand plusieurs structures respectent l'octet, on les départage par leurs charges formelles. Le cas
d'école est l'ion \textbf{perchlorate} \ce{ClO4-}.

\begin{exbox}[\ce{ClO4-} : l'octet étendu annule les charges]
\textbf{Structure A} — quatre liaisons \ce{Cl-O} \emph{simples}. Le chlore respecte l'octet mais :
\[ \CF(\ce{Cl}) = 7 - 0 - 4 = +3, \qquad \CF(\ce{O}_{\text{simple}}) = 6 - 6 - 1 = -1\ \text{(×4)}. \]
Somme $= +3 + 4\times(-1) = -1$ \checkmark, mais une charge $+3$ sur le chlore est \emph{défavorable}.

\smallskip
\textbf{Structure B} — le \ce{Cl} (3\textsuperscript{e} période) fait un \textbf{octet étendu} :
trois doubles liaisons $+$ une simple.
\[ \CF(\ce{Cl}) = 7 - 0 - 7 = 0, \qquad \CF(\ce{O}_{\text{double}}) = 6 - 4 - 2 = 0, \qquad
   \CF(\ce{O}_{\text{simple}}) = -1. \]
Charges formelles \emph{minimales} : la structure B est la plus représentative.
\end{exbox}

\begin{attnbox}[qui a le droit d'annuler ses charges ?]
Seul un atome de \textbf{période $\geq 3$} (\ce{P}, \ce{S}, \ce{Cl}\ldots) peut faire un octet
étendu pour abaisser sa charge formelle. Un atome de 2\textsuperscript{e} période (\ce{C}, \ce{N},
\ce{O}) ne peut \textbf{jamais} dépasser $8$ électrons : il garde sa charge formelle, même élevée
(ainsi \ce{NO3-} garde \ce{N} à $\CF = +1$, quand \ce{ClO4-} ramène \ce{Cl} à $0$).
\end{attnbox}

\begin{astucebox}[le réflexe en une ligne]
$\CF = N_v - 2\times(\text{doublets libres}) - (\text{nb de liaisons})$, et \textbf{une double
liaison compte pour 2 liaisons}. Vérifie toujours : la somme des \CF{} doit redonner la charge de
l'édifice.
\end{astucebox}

\end{document}
